\documentclass{article}
\usepackage{tcee-repaso}
\tcetema[Globalización y movimientos de factores]{3.B.18}{La nueva globalización económica y financiera. Determinantes y efectos de los movimientos internacionales de factores productivos: movilidad de trabajadores, inversión de cartera e inversión directa, con especial referencia a las cadenas globales de valor.}
\tcefecha{13/11/2026}
\begin{document}
\cabecerarepaso
\begin{columnas}

% ==========================================================================
\bloque{Esquema del cante}
\esquemacante{%
\apartadocante{Introducción: «el mundo es un pañuelo»}{2,5}
\apartadocante{I. La nueva globalización}{5}
\subcante{1. Concepto y medición}{1,5}
\subcante{2. Las oleadas}{1}
\subcante{3. Rasgos (5)}{2,5}
\apartadocante{II. Movimientos de factores}{13}
\subcante{1. Marco común: \autor{MacDougall}}{2}
\subcante{2. Trabajadores}{4}
\subcante{3. Inversión de cartera}{3}
\subcante{4. Inversión directa}{4}
\apartadocante{III. Cadenas globales de valor}{8}
\subcante{1. Concepto y medición}{2}
\subcante{2. Teoría de la fragmentación}{2,5}
\subcante{3. Determinantes y efectos}{2}
\subcante{4. Riesgos y reconfiguración}{1,5}
\apartadocante{Conclusión}{1,5}}
\emph{Pizarra}: intro sin pizarra, con «26,6~\%» y «3,7~\%» en una esquina. A la izquierda, los 3 bloques y el argumento común («cada factor va a donde rinde más, frenado por lo que cuesta moverse»). I: boceto de los 3 indicadores, columna de oleadas con los 3 desagregamientos, 5 palabras de los rasgos. II: \autor{MacDougall} por capas (áreas A-E) con la igualdad de salarios en morado; tabla de los 3 factores (determinantes y efectos) que se va llenando; HMY con las 3 rectas y los umbrales en morado. III: los 3 países con hacia atrás y hacia delante y el índice en morado; la tabla de 4 casillas de hacer o comprar; curva de la sonrisa.

% ==========================================================================
\bloque[2,5]{Introducción}
\begin{itemize}
\item \textbf{Enganche}: «El mundo es un pañuelo»: \dato{26,6~\%}{OCDE, 2022; 23,5~\% en 2010} del valor de lo que España exporta es valor añadido extranjero, y casi una quinta parte de su valor viaja en exportaciones ajenas. Más pequeño para las cosas que para las personas: solo el \dato{3,7~\%}{ONU} vive fuera de su país. Y se dobla: aranceles desde 2018, pandemia, guerras, «entre amigos». ¿Se acaba la globalización o cambia de forma?
\item \textbf{Relevancia}: comercio \dato{57~\% del PIB mundial}{BM, 2024}, casi la mitad en CGV. IED \dato{1,6 billones de \$}{UNCTAD, 2025}; remesas \dato{$\approx$700.000 M\$}{BM}, más que la IED que reciben los países de renta baja y media; la cartera transmite en horas a la Fed \vertema{3.B.15}. Reparto desigual $\Rightarrow$ rechazo (\autor{Rodrik}, 2011).
\item \textbf{Contextualización}: \concepto{globalización}: integración creciente de mercados de bienes, servicios, capitales, trabajo y conocimiento. \concepto{Nueva globalización} (desde $\approx$1990): las TIC abaratan coordinar a distancia; la producción se parte en tareas, las \concepto{CGV}; se multiplican los flujos brutos (\autor{Baldwin}, 2016). \concepto{Movimientos de factores}: trabajo (migración) y capital (cartera, directa). Argumento común: cada uno va a donde rinde más ($w$, $r$, $B$) y se mueve menos de lo predicho porque moverse cuesta. \emph{Convención}: IED = control o influencia ($\geq$10~\% de los votos); cartera, sin control (MBP6) \vertema{3.B.11}.
\item \textbf{Problemática}: ¿qué hay de nuevo? ¿Retrocede? ¿Por qué se mueven (y tan poco) los factores? ¿Quién gana y quién pierde? ¿Qué cambia cuando se mueven tareas?
\item \textbf{Estructura}: del fenómeno a sus causas y de los factores a las tareas: \hyperref[rep:I]{I}, nueva globalización; \hyperref[rep:II]{II}, los 3 factores con el mismo patrón; \hyperref[rep:III]{III}, CGV.
\end{itemize}

% ==========================================================================
\bloque[5]{I. La nueva globalización económica y financiera\label{rep:I}}
\begin{ideaclave}
Integra más los bienes, por tareas, y el capital, en bruto, que el trabajo. Desde 2008 se frena y desde 2018 se reordena por bloques, pero no retrocede.
\end{ideaclave}

\sub{I.1. Concepto y medición\label{rep:I.1}}
\begin{itemize}
\item 4 dimensiones: comercio, finanzas, migraciones, información.
\item \emph{Comercio}: \dato{26~\% (1970) $\to$ 57~\% (2024), máx.\ 62~\% en 2022}{BM}. Las cifras brutas cuentan varias veces: \concepto{comercio en valor añadido} \vertema{3.A.31}. Servicios: 4 modos del AGCS (el 3.º, presencia comercial, es IED).
\item \emph{Finanzas}: \concepto{netos} (saldo de la cuenta financiera), \concepto{brutos} \dato{22~\% del PIB (2007) $\to$ 6~\% (2024)}{FMI, cálculo propio}, \emph{stocks} (PII) estables desde 2007; IED, 5~\% (2007) $\to$ $\approx$1~\%.
\item \emph{Migraciones}: \dato{304 M, 3,7~\%}{ONU, 2024} (2,9~\% en 1990): el factor menos globalizado.
\item \emph{Índices}: KOF; DHL, en máximos en 2025 pese a los aranceles.
\end{itemize}
\graficorepaso[\linewidth]{indicadores-globalizacion}{Comercio, flujos brutos e IED, y migrantes, 1970-2024}

\sub{I.2. Las oleadas\label{rep:I.2}}
\begin{itemize}
\item \autor{Baldwin} (2016): 3 \concepto{desagregamientos} (\emph{unbundlings}): transporte (1820), coordinación (1990), telepresencia (ahora).
\item \emph{1870-1914}: patrón oro, 60 M de europeos a América; convergen salarios y rentas de la tierra (\autor{Heckscher}-\autor{Ohlin}; \autor{O'Rourke}-\autor{Williamson}) \vertema{3.B.20}. \emph{1914-1945}: desglobalización. \emph{1945-1973}: GATT, capital controlado, contenedor. \emph{1973-2008}: liberalización e \concepto{hiperglobalización} (China en la OMC, 2001); movilidad del capital en U (\autor{Obstfeld}-\autor{Taylor}). \emph{Desde 2008}: \emph{slowbalisation}; desde 2018, bloques.
\end{itemize}

\sub{I.3. Rasgos de la nueva globalización\label{rep:I.3}}
\begin{itemize}
\item \textbf{a) CGV y servicios}: comercio de tareas; \concepto{servicificación} de las manufacturas.
\item \textbf{b) Globalización financiera}: brutos $\gg$ netos (riesgos en los brutos, \autor{Lane}-\autor{Milesi-Ferretti}); banca global en dólares (\autor{Bruno}-\autor{Shin}); otra inversión 10~\% del PIB en 2007, negativa en 2008-2009; papel del dólar (\autor{Gopinath}); \concepto{ciclo financiero global}: para \autor{Hélène Rey}, dilema, no trilema \vertema{3.A.42}.
\item \textbf{c) Emergentes}: de proveedores a fábricas; China, 1.er exportador; desigualdad \emph{entre} países ↓ desde los 90.
\item \textbf{d) Desigualdad dentro}: \concepto{curva del elefante} (1988-2008): centro +70~\%, 1~\% más rico +65~\%, percentil 80-90 +5~\% (\autor{Lakner}-\autor{Milanovic}, 2016); 2008-2018 sin trompa; Gini mundial 70 $\to$ 60. «Shock de China»: $\approx$¼ de la caída del empleo manufacturero de EE.\,UU. (1990-2007) \vertema{3.B.5}.
\item \textbf{e) Rechazo y fragmentación}: 3 trilemas (monetario; financiero, \autor{Schoenmaker}; político, \autor{Rodrik}). Guerra comercial (2018); arancel efectivo de EE.\,UU. \dato{2,4 $\to$ 7,6~\% $\to$ 6,9~\%}{2024, 2025, ago.\ 2026}; Supremo anula los de la IEEPA (feb.\ 2026). \emph{Friend-shoring}: comercio entre bloques \dato{$-12$~\%}{FMI}, IED $-20$~\% más que dentro; conectores (México, Vietnam); coste hasta 7~\% del PIB mundial. ¿Desglobalización? No: reordenación.
\end{itemize}
\pareja{flujos-capital-brutos}{Entradas brutas por tipo, 1980-2024}{curva-elefante}{La curva del elefante, 1988-2008}

% ==========================================================================
\bloque[13]{II. Movimientos internacionales de factores\label{rep:II}}
\begin{ideaclave}
Cada factor va a donde rinde más; su movilidad aumenta la producción mundial, pero redistribuye: gana el que se mueve y el que lo complementa en destino, y pierde el que compite con él.
\end{ideaclave}

\sub{II.1. Marco común: \autor{MacDougall} (1960)\label{rep:II.1}}
\begin{itemize}
\item \textbf{a) Idea}: mover un factor a donde rinde más aumenta la producción mundial; ¿quién gana y quién pierde? Modelo de referencia de la inmigración (\autor{Borjas}).
\item \textbf{b) Modelo}. \emph{Supuestos}: Norte y Sur, un bien, rendimientos constantes, $\mathrm{PMg}_L$ decreciente; capital fijo, $\overline L$ dado; $w=\mathrm{PMg}_L$; el Norte con más capital por trabajador; emigrar no cuesta. \emph{Desarrollo}: emigran hasta igualar salarios:
\end{itemize}
\funcionclave{\mathrm{PMg}_{L,N}(L_N)=\mathrm{PMg}_{L,S}(\overline L-L_N)=w^1}
\begin{itemize}
\item Ganancia mundial B + C. Norte: nativos pierden A, capital gana A + B; \concepto{excedente de la inmigración} B, pequeño ($\approx$0,1~\% del PIB, \autor{Borjas}, 1995). Sur: trabajadores ganan D, capital pierde D + E. Emigrantes ganan C + E.
\item \textbf{iii) Implicaciones}: eficiente, no Pareto sin compensaciones \vertema{3.A.22}; con el capital en el eje, la IED.
\item \textbf{c) Extensiones}: costes de moverse; comercio y factores \emph{sustitutivos} (\autor{Mundell}, 1957, en \autor{Heckscher}-\autor{Ohlin}) o \emph{complementarios} (\autor{Markusen}, 1983: CGV).
\item \textbf{d) Evidencia}: convergencia de 1870-1914; hoy, salarios muy distintos.
\item \textbf{e) Valoración}: ganancia potencial del \dato{67-147~\% del PIB mundial}{\autor{Clemens}, 2011}; pero estático (trabajo homogéneo, capital fijo).
\end{itemize}
\graficorepaso[0.85\linewidth]{macdougall}{Igualdad de salarios, ganancia B + C y reparto}

\sub{II.2. Movilidad de trabajadores\label{rep:II.2}}
\begin{itemize}
\item \dato{304 M}{2024}; EE.\,UU. 52 M; España, \dato{10,3 M de 49,8}{INE, 2026}, 1 de cada 5 \vertema{4.A.2}.
\item \textbf{a) Determinantes}. \emph{Decisión}: diferencia de salarios (\autor{Hicks}, 1932); inversión en capital humano (\autor{Sjaastad}, 1962):
\end{itemize}
\begin{formula}
\[\sum_{t=1}^{T}\frac{w_{N,t}-w_{S,t}}{(1+\rho)^t}>c_m\]
\end{formula}
\begin{itemize}
\item Jóvenes, destinos cercanos o con compatriotas. \emph{Selección} (\autor{Borjas}, 1987; \autor{Roy}, 1951): positiva si la cualificación se paga más en destino. \emph{Agregados}: gravedad ($>$ la mitad de los flujos, FMI): redes, distancia, lengua, conflicto; «joroba» del desarrollo (\autor{Clemens}, 2014).
\item \textbf{b) Efectos en destino}: el capital se ajusta; sustitutivos y complementarios: $-3/-4$~\% por +10~\% de oferta del grupo (\autor{Borjas}, 2003) frente a efecto casi nulo (\autor{Ottaviano}-\autor{Peri}, 2012); tareas de comunicación (\autor{Peri}-\autor{Sparber}); patentes +9-18~\%; fiscal $\pm$0,5~\% del PIB (OCDE); producción +1~\% a 5 años (FMI); ajuste en unión monetaria \vertema{3.B.16}.
\item \emph{El Mariel} (1980): 125.000 cubanos, +7~\% de la población activa de Miami; sin efectos (\autor{Card}, 1990; Nobel 2021, experimento natural \vertema{4.B.6}); \autor{Borjas} (2017) sí; \autor{Peri}-\autor{Yasenov} (2019), control sintético, no; \autor{Clemens}-\autor{Hunt}: composición de la muestra. Reacción política (\autor{Tabellini}, 2020).
\item \emph{En origen}: fuga y ganancia de cerebros (\autor{Beine}-\autor{Docquier}-\autor{Rapoport}); \concepto{remesas} \dato{$\approx$700.000 M\$}{2025}, > IED desde 2022; Tayikistán 47~\%, Tonga 39~\%, Nicaragua 27~\% del PIB; estables y anticíclicas; riesgo de enfermedad holandesa; renta secundaria \vertema{3.B.11}.
\end{itemize}
\pareja{remesas-ied-aod}{Remesas, IED y AOD, 1990-2025}{remesas-pib-paises}{Remesas en \% del PIB, 2024}

\sub{II.3. Inversión de cartera\label{rep:II.3}}
\begin{itemize}
\item \textbf{a) Determinantes}. \emph{Rentabilidad}: enfoque intertemporal ($r^A>r^*$ $\Rightarrow$ importa capital) \vertema{3.B.12}; solo netos. Puzles: \autor{Feldstein}-\autor{Horioka} (1980), \autor{Lucas} (1990: de 58 veces a 5), asignación (\autor{Gourinchas}-\autor{Jeanne}).
\item \emph{Riesgo}: diversificación (\autor{Markowitz}, 1952; \autor{Grubel}, 1968; \autor{Solnik}, 1974); varianza a medias $\tfrac{\sigma^2}{2}[1+\mathrm{corr}]$. \concepto{Sesgo nacional} (\emph{home bias}): > 90~\% en casa en EE.\,UU. y Japón (\autor{French}-\autor{Poterba}, 1991); explicaciones: costes, cobertura, gobierno corporativo, familiaridad.
\item \emph{Push y pull} (\autor{Calvo}-\autor{Leiderman}-\autor{Reinhart}, 1993); mandan los globales (\autor{Forbes}-\autor{Warnock}, 2012): «el dinero es cobarde».
\item \textbf{b) Efectos}: beneficios teóricos (asignación, riesgo, disciplina) difíciles de encontrar (\autor{Kose} et al., 2009; consumo menos correlacionado que la producción); mala asignación en el sur de Europa (\autor{Gopinath} et al., 2017); paradas súbitas \vertema{3.B.17}; la IED, la más estable; gestión de flujos aceptada.
\end{itemize}

\sub{II.4. Inversión directa\label{rep:II.4}}
\begin{itemize}
\item \textbf{a) Determinantes}: IED $\geq$10~\%; \emph{greenfield} o fusiones; \concepto{horizontal}, \concepto{vertical}, \concepto{plataforma}. \dato{1,6 billones, +6~\%}{UNCTAD, 2025}; 20 primeros > 80~\%; sectores estratégicos 44~\% (16~\% en 2020); sobre todo rico-rico y horizontal.
\item \textbf{b) Teorías}: ventajas propias (\autor{Hymer}); internalización (\autor{Buckley}-\autor{Casson}) \vertema{3.A.15}; \concepto{OLI} (\autor{Dunning}): propiedad, localización, internalización; ciclo del producto (\autor{Vernon}, 1966); vertical (\autor{Helpman}, 1984); capital-conocimiento (\autor{Markusen}).
\item \textbf{HMY (2004)}. \emph{Idea}: \concepto{disyuntiva proximidad-concentración}. \emph{Supuestos}: \autor{Melitz} (CES, $\sigma>1$, productividad $\varphi$); exportar ($wF^x$, $\tau$) o filial ($wF^I$ mayor, sin transporte). \emph{Desarrollo}: rectas en $\varphi^{\sigma-1}$; $B^x$ con pendiente $\tau^{1-\sigma}$ veces la de $B^I$:
\end{itemize}
\funcionclave{F^{I}>\tau^{\sigma-1}F^x>F^d\Rightarrow\varphi^*<\varphi^*_x<\varphi^*_{I}}
\graficorepaso[0.8\linewidth]{hmy}{Exportar o invertir: las 3 rectas y los umbrales}
\begin{itemize}
\item Cierran, solo nacional, exportan, filial: «ver los toros desde la barrera» es exportar; bajar al ruedo, invertir. \emph{Implicaciones}: más IED con costes de comerciar altos, costes fijos de planta bajos y más dispersión; \emph{tariff jumping} \vertema{3.B.10}. \emph{Evidencia}: multinacionales +15~\% más productivas. \emph{Valoración}: deja fuera la vertical y hacer o comprar.
\item \textbf{c) Efectos}: destino: capital, empleo; desbordamientos horizontales pequeños (\autor{Aitken}-\autor{Harrison}), verticales positivos (\autor{Javorcik}); traslado de beneficios \vertema{4.B.18}. Origen: complementarios (\autor{Desai}-\autor{Foley}-\autor{Hines}); prima de cualificación ↑ en los 2 (\autor{Feenstra}-\autor{Hanson}, 1997).
\end{itemize}

% ==========================================================================
\bloque[8]{III. Cadenas globales de valor\label{rep:III}}
\begin{ideaclave}
Se comercia con tareas: cada país hace una fase, y la empresa decide dónde y si la hace ella. Dan crecimiento a quien se integra, reparten el valor hacia los extremos y se gestionan mejor diversificando que deshaciéndolas.
\end{ideaclave}

\sub{III.1. Concepto y medición\label{rep:III.1}}
\begin{itemize}
\item \concepto{CGV}: fases de un producto con al menos 2 en países distintos. Causas: contenedor, TIC (2.º desagregamiento), liberalización y acuerdos profundos. Casi la mitad del comercio; estancadas desde 2008.
\item Medición: \concepto{tablas \emph{input-output} multirregionales} (WIOD, TiVA, FIGARO) e inversa de \autor{Leontief}. \concepto{Hacia atrás} (VAE) y \concepto{hacia delante} (VAI); \concepto{distancia a la demanda final} (\autor{Antràs}, 2012):
\end{itemize}
\funcionclave{\text{participación}=\frac{\mathrm{VAE}}{X}+\frac{\mathrm{VAI}}{X}}
\begin{itemize}
\item España \dato{26,6~\%}{OCDE, 2022}; OCDE 30,4~\%, EE.\,UU. 8,3~\%, China 17~\%.
\end{itemize}
\graficorepaso[0.9\linewidth]{cgv-participacion}{Hacia atrás y hacia delante (figura)}

\sub{III.2. Teoría de la fragmentación\label{rep:III.2}}
\begin{itemize}
\item \emph{Comercio de tareas} (\autor{Grossman}-\autor{Rossi-Hansberg}, 2008): efecto productividad, que puede subir el salario de los poco cualificados que se quedan.
\item \emph{Hacer o comprar entre países}: contratos incompletos (\autor{Grossman}-\autor{Hart}, 1986; Nobel de \autor{Hart}, 2016); se integran las fases donde pesa la inversión de la matriz (\autor{Antràs}, 2003); las más productivas integran fuera, luego subcontratan fuera, luego en casa (\autor{Antràs}-\autor{Helpman}, 2004). Cadenas secuenciales (\autor{Antràs}-\autor{Chor}, 2013); geografía (\autor{Antràs}-\autor{de Gortari}, 2020).
\end{itemize}
\begin{formula}
\centering\small
\begin{tabular}{@{}p{0.2\linewidth}>{\centering\arraybackslash}p{0.36\linewidth}>{\centering\arraybackslash}p{0.36\linewidth}@{}}
& \textbf{Integración} & \textbf{Subcontratación}\\
En el país & nacional & nacional\\
En otro país & \textbf{IED vertical (intrafirma)} & \emph{offshore outsourcing}\\
\end{tabular}
\end{formula}

\sub{III.3. Determinantes y efectos\label{rep:III.3}}
\begin{itemize}
\item \textbf{a) Determinantes} (BM, 2020): dotaciones, tamaño, geografía («fábricas» Europa, América del Norte, Asia), instituciones, políticas \vertema{3.B.10}. 4 posiciones: materias primas, manufactura limitada, avanzada y servicios, innovadoras (España entró en estas entre 1990 y 2015).
\item \textbf{b) Efectos}: +1~\% de participación $\Rightarrow$ > +1~\% de renta pc; industrializarse sin hacer la cadena entera; \concepto{curva de la sonrisa} (\autor{Stan Shih}): el valor en los extremos; transmisión de perturbaciones, aranceles sobre insumos \vertema{3.B.7}; huella de carbono \vertema{3.B.31}.
\end{itemize}
\graficorepaso[0.6\linewidth]{curva-sonrisa}{La curva de la sonrisa}

\sub{III.4. Riesgos y reconfiguración\label{rep:III.4}}
\begin{itemize}
\item Choques: pandemia, chips, Suez, Ucrania, mar Rojo. Respuestas: \emph{reshoring}, \emph{nearshoring}, \emph{friend-shoring}; IRA y CHIPS (2022), ley europea de chips (2023) \vertema{3.B.41}.
\item Datos: la cadena da un rodeo (Vietnam, México: \autor{Freund}; \autor{Alfaro}-\autor{Chor}); relocalizar no da resiliencia (\autor{Bonadio} et al., 2021); diversificar (\autor{Baldwin}-\autor{Freeman}). Reglas en acuerdos profundos \vertema{3.B.33}.
\item España: \dato{25,7~\% hacia atrás, 19,9~\% hacia delante}{Eurostat FIGARO, 2024}; UE, 38,4~\% \vertema{4.A.22}.
\end{itemize}
\graficorepaso[0.9\linewidth]{contenido-importado-ue}{Participación en las CGV de la UE, 2024}

% ==========================================================================
\bloque[1,5]{Conclusión}
\begin{itemize}
\item \textbf{Recapitulación}: bienes por tareas y capital en bruto, más que trabajo; frenazo (2008) y bloques (2018). Los factores van a donde rinden más: trabajo poco, cartera mucho y con miedo, IED de las mejores empresas. CGV: motor de la integración de los emergentes.
\item \textbf{Relevancia}: ganancias grandes y repartidas, costes concentrados: rechazo y fragmentación.
\item \textbf{Extensiones}: \vertema{3.B.12}, \vertema{3.B.15}, \vertema{3.B.17}, \vertema{3.B.19}, \vertema{3.B.39}; España: \vertema{4.A.22}, \vertema{4.A.23}, \vertema{4.B.18}.
\item \textbf{Opinión}: resiliencia no es autarquía; compensar a quien pierde mejor que levantar barreras.
\item \textbf{Idea final}: el mundo sigue siendo un pañuelo, pero se dobla por las costuras de los bloques: la pregunta es cómo repartir lo que se gana.
\end{itemize}

% ==========================================================================
\bloque{Autores y fechas\label{rep:autores}}
\hito{1932}{\autor{Hicks}: salarios y migración}
\hito{1957}{\autor{Mundell}: comercio y factores, sustitutivos}
\hito{1960}{\autor{MacDougall}: inversión extranjera}
\hito{1962}{\autor{Sjaastad}: emigrar es invertir}
\hito{1966}{\autor{Vernon}: ciclo del producto}
\hito{1977}{\autor{Dunning}: paradigma OLI}
\hito{1987}{\autor{Borjas}: selección}
\hito{1990}{\autor{Card}: el Mariel}
\hito{1993}{\autor{Calvo}-\autor{Leiderman}-\autor{Reinhart}: \emph{push} y \emph{pull}}
\hito{2001}{China en la OMC}
\hito{2003}{\autor{Antràs}: hacer o comprar}
\hito{2004}{\autor{Helpman}-\autor{Melitz}-\autor{Yeaple}}
\hito{2008}{\autor{Grossman}-\autor{Rossi-Hansberg}: tareas}
\hito{2013}{\autor{Rey}: ciclo financiero global}
\hito{2016}{\autor{Baldwin}: 3 desagregamientos; elefante}
\hito{2018}{Guerra comercial EE.\,UU.-China}

% ==========================================================================
\bloque{Datos clave}
\begin{itemize}
\item Comercio 57~\% del PIB mundial (2024; 26~\% en 1970). Flujos brutos 6~\% (22~\% en 2007). IED 1,6 billones de \$ (2025). Migrantes 304 M, 3,7~\%. Remesas $\approx$700.000 M\$. España: 10,3 M de nacidos fuera; contenido extranjero de sus exportaciones 26,6~\% (OCDE, 2022) o 25,7~\% (FIGARO, 2024). Arancel efectivo de EE.\,UU.: 2,4 $\to$ 7,6~\% (2024-2025).
\end{itemize}

% ==========================================================================
\bloque{Preguntas del tribunal}
\begin{enumerate}
\item \textbf{¿Quién pierde con la inmigración?} En \autor{MacDougall}, los nativos sustitutivos (A) a corto; el excedente B es pequeño.
\item \textbf{¿Comercio y factores, sustitutivos o complementarios?} Con dotaciones, sustitutivos (\autor{Mundell}); con tecnología o escala, complementarios (\autor{Markusen}).
\item \textbf{¿Por qué no va el capital a los pobres?} \autor{Lucas}: capital humano, instituciones y riesgo.
\item \textbf{¿Qué empresas abren filiales?} Las más productivas (HMY).
\item \textbf{¿Hay desglobalización?} No: frenazo y reordenación por bloques; el comercio da un rodeo.
\end{enumerate}

\begin{trampa}
\begin{itemize}
\item IED: $\geq$10~\% de los votos; cartera, por debajo, y deuda.
\item Hacia atrás = VAE (extranjero en mis exportaciones); hacia delante = VAI (mío en las de otros).
\item Remesas: renta secundaria, no cuenta financiera.
\item El modo 3 del AGCS (presencia comercial) es IED.
\item Remesas > IED en los países de renta baja y media desde 2022, no antes.
\item Curva del elefante: los que menos ganan son el percentil 80-90 (clases medias y bajas de los ricos), no los más pobres.
\item \autor{MacDougall}: el excedente de la inmigración (B) es pequeño; la redistribución (A), grande.
\item Sesgo nacional: se tiene \emph{menos} activo extranjero del que aconseja la teoría.
\end{itemize}
\end{trampa}

\end{columnas}
\end{document}
